\documentclass[11pt]{article}
\usepackage[a4paper,margin=1in]{geometry}
\usepackage{amsmath,amssymb,amsthm,microtype}
\newtheorem{theorem}{Theorem}[section]
\theoremstyle{remark}\newtheorem{remark}[theorem]{Remark}
\title{Parity-Orthogonal Hardy Channels at the Endpoint Logarithm}
\author{Bin Wang}\date{July 2026}
\begin{document}\maketitle
\begin{abstract}
The endpoint singularity found in RH-312 admits an exact even--odd
decomposition.  We identify the two closed forms and prove that the Hardy
norm is their orthogonal sum.  Endpoint convergence is therefore equivalent
to convergence in both parity channels separately.  This criterion does not
assert that either actual noisy channel converges.
\end{abstract}

\section{The two singular channels}
Let $L(w)=\log(1-w)+w=-\sum_{n\ge2}w^n/n$.  Its parity projections are
\[
 L_{\rm ev}(w)=\frac{L(w)+L(-w)}2=\frac12\log(1-w^2),
\]
and
\[
 L_{\rm odd}(w)=\frac{L(w)-L(-w)}2
 =w+\frac12\log\frac{1-w}{1+w}.
\]
Both belong to $H^2$ on the unit disk, and their Taylor supports are disjoint.

\begin{theorem}[orthogonal endpoint split]
For every $f\in H^2$, define
$f_{\rm ev}(w)=\{f(w)+f(-w)\}/2$ and
$f_{\rm odd}(w)=\{f(w)-f(-w)\}/2$.  Then
\[
 f=f_{\rm ev}+f_{\rm odd},\qquad
 \|f\|_{H^2}^2=\|f_{\rm ev}\|_{H^2}^2+
 \|f_{\rm odd}\|_{H^2}^2.
\]
Consequently $f_\sigma\to0$ in $H^2$ if and only if both parity projections
tend to zero.
\end{theorem}
\begin{proof}
The even and odd Taylor coefficients occupy disjoint subsets of the
orthonormal monomial basis.  Parseval gives the identity and the equivalence.
\end{proof}

\section{Application to the deterministic endpoint}
RH-312 writes the endpoint target as
$L+H_{\rm reg}$ with $H_{\rm reg}$ analytic past the unit circle.  Projecting
gives
\[
 (L_{\rm ev}+H_{{\rm reg},{\rm ev}})
 +(L_{\rm odd}+H_{{\rm reg},{\rm odd}}),
\]
and both regular pieces retain the larger analytic radius.  Thus the two
displayed logarithms contain the entire parity-resolved endpoint
singularity.

\begin{remark}
The theorem is an exact criterion, not a proof that the noisy modulus
complement approximates either channel.  No spectral cloud, Gate A--E, or RH
conclusion is obtained.
\end{remark}
\end{document}
